\section{A consequence for unions of two sets}\label{sec:consequence}

The comparison also strengthens a recent bound for sets that cannot be
covered by two members of a large class. For $\mathcal A\subseteq2^X$,
write $\mu_p(\mathcal A)=\P(X_p\in\mathcal A)$ and define
\[
 \operatorname{Bad}_2(\mathcal A)
 =\{S\subseteq X:S\nsubseteq A_1\cup A_2
                   \text{ for all }A_1,A_2\in\mathcal A\}.
\]
When $\mu_p(\mathcal A)>1/2$, Li's fractional discrete-convexity
theorem~\cite[Corollary~3.1]{Li2026} provides a fractional cover for this
family. Combining that input with Theorem~\ref{thm:main} gives the
following integral conclusion.

\begin{corollary}\label{cor:two-unions}
Let $X$ be a finite nonempty set, let $0<p<1$, and let
$\mathcal A\subseteq2^X$ satisfy $\mu_p(\mathcal A)>1/2$.
Set $C=25\cdot512^4$. There is a family $\mathcal G\subseteq2^X$ such
that
\[
 \operatorname{Bad}_2(\mathcal A)\subseteq\langle\mathcal G\rangle,
 \qquad
 \sum_{S\in\mathcal G}\left(\frac{p}{2C}\right)^{|S|}\le\frac12.
\]
In particular, $\operatorname{Bad}_2(\mathcal A)$ is
$(p/(2C))$-small. No monotonicity assumption on $\mathcal A$ is needed.
\end{corollary}

\begin{proof}
If $\operatorname{Bad}_2(\mathcal A)$ is empty, take
$\mathcal G=\varnothing$. Otherwise this is a nonempty increasing
family. It is proper because $\mathcal A$ is nonempty and hence
$\varnothing\notin\operatorname{Bad}_2(\mathcal A)$.
Li's theorem, in the covering formulation given by
Park~\cite[Theorem~3.4]{Park2026}, supplies a fractional cover of
$\operatorname{Bad}_2(\mathcal A)$ of cost at most $1/2$ at parameter
$p/2$. Its coefficients may be capped at one without affecting coverage
or increasing cost. Theorem~\ref{thm:main} therefore yields an integral
cover of cost at most $1/2$ at parameter $p/(2C)$.
\end{proof}

Park~\cite[Corollary~3.5]{Park2026} obtained the same conclusion with
smallness parameter
\[
 \frac{p}{K\max\{1,\log\log(1/p)\}}
\]
for a universal $K$. Corollary~\ref{cor:two-unions} removes the double
logarithmic loss
while retaining unions of two members and the covering budget $1/2$.
